% TODO checklist:
% - [x] Inspected src/logging_setup.py (structlog configuration)
% - [x] Inspected src/observability/middleware.py (request context binding)
% - [x] Reviewed Q&A.md sections on logging and structured logging
% - [x] Reviewed README.md for logging overview
% - [x] Examined AccessPathFilter implementation for noise reduction

\section{Logging}
\label{sec:logging}

The Cineca Agentic Platform implements a sophisticated structured logging framework built on \texttt{structlog}, providing machine-parseable JSON logs for production environments and human-readable console output for development. The logging subsystem is designed to support distributed tracing, request correlation, and efficient log aggregation.

\subsection{Logging Architecture Overview}

The logging architecture follows a unified pipeline pattern that bridges Python's standard \texttt{logging} module, FastAPI/Uvicorn loggers, and \texttt{structlog} processors into a single coherent output stream. This design ensures that all log entries---regardless of their origin---pass through the same processing chain and emit consistent, structured output.

\begin{lstlisting}[language=Python, caption={Logging pipeline initialization (src/logging\_setup.py)}]
def setup_logging(level: str | int = "INFO") -> None:
    """
    Initialize structlog + stdlib logging with a single StreamHandler.
    Idempotent: calling again replaces handlers with current configuration.
    """
    log_level = _coerce_level(level)
    use_json = _wants_json()

    # Base processors for timestamp and context
    timestamper = structlog.processors.TimeStamper(fmt="iso", utc=True)
    pre_chain = [
        structlog.contextvars.merge_contextvars,
        structlog.stdlib.add_log_level,
        structlog.stdlib.add_logger_name,
        timestamper,
    ]

    # Renderer selection based on environment
    if use_json:
        renderer = structlog.processors.JSONRenderer(sort_keys=True)
    else:
        renderer = structlog.dev.ConsoleRenderer()
\end{lstlisting}

\subsection{Environment-Aware Log Formatting}

The logging system automatically selects the appropriate output format based on the deployment environment:

\begin{itemize}
    \item \textbf{Development Mode}: Uses \texttt{structlog.dev.ConsoleRenderer()} which produces colorized, human-readable output with key-value pairs aligned for easy scanning.
    
    \item \textbf{Production Mode}: Uses \texttt{structlog.processors.JSONRenderer()} which emits single-line JSON objects suitable for ingestion by log aggregation systems like Elasticsearch, Splunk, or CloudWatch.
\end{itemize}

The format selection is determined by the \texttt{LOG\_FORMAT} environment variable or derived from \texttt{APP\_ENV}:

\begin{lstlisting}[language=Python, caption={Log format selection logic}]
def _wants_json() -> bool:
    fmt = os.getenv("LOG_FORMAT", "").strip().lower()
    if fmt in {"json", "console"}:
        return fmt == "json"
    env = os.getenv("APP_ENV", "dev").strip().lower()
    return env in {"prod", "production"}
\end{lstlisting}

\subsection{Structured Log Schema}

Production JSON logs follow a consistent schema that supports efficient querying and correlation:

\begin{lstlisting}[language=json, caption={Example structured log entry}]
{
  "timestamp": "2025-01-15T10:30:45.123456Z",
  "level": "info",
  "logger": "src.routers.agents",
  "event": "http.request",
  "request_id": "a1b2c3d4e5f6",
  "trace_id": "00000000000000001234567890abcdef",
  "method": "POST",
  "path": "/v1/agent-runs",
  "status": 200,
  "duration_s": 0.342156,
  "user_agent": "Mozilla/5.0...",
  "client_ip": "10.0.0.1"
}
\end{lstlisting}

Key fields in the log schema include:

\begin{table}[ht]
\centering
\caption{Standard structured log fields.}
\label{tab:log-fields}
\begin{tabularx}{\textwidth}{llX}
\hline
\textbf{Field} & \textbf{Type} & \textbf{Description} \\
\hline
\texttt{timestamp} & ISO 8601 & UTC timestamp in ISO format with microsecond precision \\
\texttt{level} & string & Log level (debug, info, warning, error, critical) \\
\texttt{logger} & string & Fully-qualified logger name \\
\texttt{event} & string & Semantic event identifier (e.g., ``http.request'', ``agent.run.complete'') \\
\texttt{request\_id} & string & Correlation ID for request tracing \\
\texttt{trace\_id} & string & OpenTelemetry trace ID (when tracing is enabled) \\
\texttt{tenant\_id} & string & Multi-tenant isolation identifier \\
\hline
\end{tabularx}
\end{table}

\subsection{Request Context Correlation}

The \texttt{ObservabilityMiddleware} binds request-specific context to all log entries generated during request processing:

\begin{lstlisting}[language=Python, caption={Request context binding in middleware}]
async def dispatch(self, request: Request, call_next: Callable) -> Response:
    # Correlation id (prefer incoming header)
    request_id = request.headers.get("x-request-id") or uuid.uuid4().hex
    request.state.request_id = request_id

    # Bind logging context
    bind_contextvars(
        request_id=request_id,
        method=request.method,
        path=_get_route_template(request),
        client_ip=request.client.host if request.client else None,
    )

    # Optionally bind OpenTelemetry trace ID
    trace_id = _current_trace_id()
    if trace_id:
        bind_contextvars(trace_id=trace_id)

    try:
        response = await call_next(request)
        # ... log completion ...
    finally:
        # Clean up per-request context
        unbind_contextvars("request_id", "method", "path", 
                          "client_ip", "trace_id")

    return response
\end{lstlisting}

This pattern ensures that all log entries generated during a request---whether from the router, service layer, repository, or external adapters---include the same correlation context for efficient debugging and tracing.

\subsection{Log Level Configuration}

The platform supports standard Python log levels with sensible defaults:

\begin{table}[ht]
\centering
\caption{Log level configuration and use cases.}
\label{tab:log-levels}
\begin{tabular}{llp{7cm}}
\hline
\textbf{Level} & \textbf{Default Use} & \textbf{Description} \\
\hline
\texttt{DEBUG} & Development & Detailed diagnostic information for debugging \\
\texttt{INFO} & Production & Normal operational events (request completion, job status) \\
\texttt{WARNING} & Production & Recoverable issues (rate limits, fallback activation) \\
\texttt{ERROR} & Production & Errors requiring attention (failed requests, exceptions) \\
\texttt{CRITICAL} & Production & System-level failures (database unavailable, security breach) \\
\hline
\end{tabular}
\end{table}

Log levels are configured through the \texttt{LOG\_LEVEL} environment variable or programmatically during application startup.

\subsection{High-Frequency Log Filtering}

To prevent log noise from overwhelming production systems, the platform implements an \texttt{AccessPathFilter} that suppresses high-frequency access logs for health and metrics endpoints:

\begin{lstlisting}[language=Python, caption={Access log noise filtering}]
class AccessPathFilter(logging.Filter):
    """Filter that drops records containing suppressed path fragments."""

    def __init__(self, *paths: str) -> None:
        super().__init__(name="access-path-filter")
        self._paths = tuple(p.lower() for p in paths if p)

    def filter(self, record: logging.LogRecord) -> bool:
        if not self._paths:
            return True
        targets = [record.getMessage().lower()]
        # ... extract additional targets from args ...
        return not any(
            path in target for target in targets for path in self._paths
        )

# Applied to uvicorn.access logger
noise_paths = ["/metrics", "/health", "get /v1/agent-runs/", 
               "options /v1/"]
access_filter = AccessPathFilter(*noise_paths)
logging.getLogger("uvicorn.access").addFilter(access_filter)
\end{lstlisting}

Additional paths can be configured via the \texttt{UVICORN\_ACCESS\_FILTER\_PATHS} environment variable (comma-separated).

\subsection{Library Noise Reduction}

The platform reduces log noise from third-party libraries by elevating their minimum log level:

\begin{lstlisting}[language=Python, caption={Library log level adjustment}]
# Reduce noise from common libraries
for noisy in ("asyncio", "httpx", "urllib3"):
    logging.getLogger(noisy).setLevel(
        min(log_level, logging.WARNING)
    )
\end{lstlisting}

This ensures that library debug and info messages don't pollute application logs while still capturing warnings and errors that may indicate underlying issues.

\subsection{Logger Unification}

The platform unifies log output from multiple sources by configuring all relevant loggers to use the same handler and formatter:

\begin{lstlisting}[language=Python, caption={Logger unification}]
# Build ProcessorFormatter for stdlib handlers
formatter = structlog.stdlib.ProcessorFormatter(
    processor=renderer,
    foreign_pre_chain=pre_chain,
)

# Create single stream handler
handler = logging.StreamHandler(sys.stdout)
handler.setFormatter(formatter)
handler.setLevel(log_level)

# Apply to root and framework loggers
root = logging.getLogger()
root.handlers[:] = [handler]
root.setLevel(log_level)

for name in ("uvicorn", "uvicorn.error", "uvicorn.access", "fastapi"):
    lg = logging.getLogger(name)
    lg.handlers[:] = [handler]
    lg.setLevel(log_level)
    lg.propagate = False  # Prevent duplicate entries
\end{lstlisting}

\subsection{Structured Logging Best Practices}

The platform's logging implementation follows several best practices for enterprise observability:

\begin{enumerate}
    \item \textbf{Semantic Event Names}: Log events use hierarchical, dot-separated names (e.g., \texttt{http.request}, \texttt{agent.run.complete}, \texttt{health.probe.failed}) that enable efficient filtering and aggregation.
    
    \item \textbf{Consistent Key Names}: All log entries use snake\_case keys consistently, avoiding abbreviations that may cause confusion.
    
    \item \textbf{Bounded Value Sizes}: Large values (request bodies, LLM responses) are truncated or summarized to prevent log storage bloat.
    
    \item \textbf{Sensitive Data Exclusion}: Passwords, tokens, and PII are never logged; the PII scrubbing subsystem (covered in Chapter~\ref{ch:security}) operates at a separate layer but complements logging hygiene.
    
    \item \textbf{Contextual Enrichment}: Logs include tenant ID, user ID, and request ID where available, enabling multi-dimensional analysis.
\end{enumerate}

\subsection{Integration with Log Aggregation Systems}

The JSON log format is designed for direct ingestion by modern log aggregation platforms:

\begin{itemize}
    \item \textbf{Elasticsearch/OpenSearch}: JSON logs can be indexed directly with automatic field type detection.
    
    \item \textbf{Grafana Loki}: Labels extracted from JSON fields enable efficient log querying via LogQL.
    
    \item \textbf{CloudWatch Logs}: AWS CloudWatch natively parses JSON logs for Insights queries.
    
    \item \textbf{Splunk}: JSON logs can be parsed with minimal configuration using Splunk's automatic field extraction.
\end{itemize}

For containerized deployments, logs are written to stdout and captured by the container runtime (Docker, Kubernetes) for forwarding to the aggregation backend.

\subsection{Audit-Specific Logging}

Beyond operational logging, the platform maintains a separate audit logging channel (covered in Section~\ref{sec:audit-logging}) for security-sensitive events that require immutable, long-term retention. The operational logging system described in this section focuses on diagnostic and performance-related events.

